% P1 学术写作版（2026-11-24）：引言、研究目的与主要贡献。
% 由 cellcharter_study_introduction_module_rev_p0_20261124.tex 改写；REV_P0/V1 保持不变。
% 每条贡献的证据来源见 P1_20261124/P1_EVIDENCE_MAP_20261124.md。未编译。
\section{引言}
\label{sec:cc-intro-p1}

\subsection{研究背景与问题}
空间转录组技术在测量基因表达的同时保留了组织中的采样位置，使研究者能够依据分子特征划分组织的空间结构域。空间域识别已成为理解组织分层、肿瘤微环境和细胞间空间关系的重要分析步骤\cite{maynard2021,wu2021}。为把空间上下文引入聚类，现有方法发展出多种机制，包括空间图表示学习\cite{xu2024,dong2022stagate}、邻域增强表示\cite{singhal2024}、空间先验统计模型\cite{zhao2021}以及多尺度图融合\cite{zhang2025}等。

CellCharter是其中一种可扩展的成熟框架\cite{varrone2024}。它先由scVI得到每个观测单元的低维表达表示\cite{lopez2018}，再把多层空间邻域的平均特征与自身特征拼接，最后用高斯混合模型聚类，并以聚类稳定性辅助选择簇数。原作者在DLPFC基准数据上把邻域层数作为超参数调优，最终选用4层邻域，并报告CellCharter在多样本联合聚类中优于对照方法\cite{varrone2024,cellcharteranalyses2024}。

对于这类已被广泛采用的方法，一个实际而尚未被充分回答的问题是：\textbf{在不同的使用与评价条件下，“更大的邻域尺度带来更好的空间域”这一结论能否成立，又在多大程度上成立？}现有基准通常在参考类别数已知的条件下，以ARI这类整体一致性指标比较方法。但在实际分析中，至少还有三个因素会影响结论：
\begin{itemize}
  \item 聚类数往往需要在不知道参考类别数的情况下确定；
  \item 整体一致性高，并不等于层间边界被准确定位；
  \item 同一流程在不同供体、不同组织或不同患者中的表现可能差异很大。
\end{itemize}
如果只报告单一ARI均值，邻域尺度的作用就可能与聚类数设定、指标定义和样本异质性混在一起。已有的空间聚类评价研究也指出，在数据处理、标签质量和参数设定不同时，方法排序具有不确定性\cite{sun2026}。

\subsection{研究目的}
本研究不提出新的聚类算法，而是把CellCharter的邻域聚合与GMM聚类流程作为研究对象，系统考察其性能结论在什么条件下成立。具体围绕以下四个问题：
\begin{enumerate}
  \item \textbf{邻域表征：}邻域聚合的深度和组成如何影响空间域划分？收益集中在哪些邻域层？这一影响在切片和供体之间是否一致？
  \item \textbf{聚类数规则：}参考类别数设定、固定K网格和无标签稳定性选K，是否会改变邻域尺度之间的比较结论？
  \item \textbf{评价指标：}整体分区一致性与空间边界定位是否给出相同的尺度排序？如果不同，差异的来源是什么？
  \item \textbf{适用范围：}与非空间及空间成熟方法相比，CellCharter处于什么位置？把同一流程用于另一组织的独立患者队列时，结论是否仍然成立？
\end{enumerate}

\subsection{研究内容与主要贡献}
围绕上述问题，本研究以3名供体的12张DLPFC Visium切片为主要数据，以6例乳腺癌Visium样本为外部队列，建立了一套多维评价设计：邻域深度与等维组成、三类聚类数条件、整体与边界两类指标、切片/供体/患者三个层级的汇总。原有方法的提出与原作者的基准结果归属CellCharter原作者；本研究独立完成的工作和得到的主要认识如下：

\begin{enumerate}
  \item \textbf{构建了面向成熟空间聚类方法的多维性能评价框架。}原作者以参考类别数下的ARI为主要指标。在此基础上，本研究把邻域层数扩展到0层和1层，加入非连续的等维邻域块组合，并在全部12张切片上设置多个GMM随机种子；同时纳入固定K网格、逐切片无标签选K以及严格与容差边界指标，使邻域尺度的效应可以在不同聚类数条件和不同评价定义下分别检验。
  \item \textbf{量化了邻域聚合收益的集中性与异质性。}第一层邻域带来的ARI提升最大，而且在全部切片上方向一致；15维之后的增益小于切片间差异，约半数切片和其中一名供体方向相反。不同的等维邻域块组合之间没有一种稳定占优；在开发切片上看似有利的远邻替换方案，在其余切片和外部队列中没有复现其优势。
  \item \textbf{揭示了聚类数规则与尺度结论的交互作用。}在固定K时，尺度排序在各K下稳定；但在多数七层切片上，ARI在少于参考层数的K处最大，说明参考类别数不一定是ARI最优的设定。逐切片的无标签稳定性选K给出的簇数离散，会改变尺度排序。
  \item \textbf{揭示了整体一致性与严格边界定位之间的评价权衡，并解释了其来源。}更深的邻域表征提高了容差边界匹配和整体一致性，严格边界F1却下降。分解precision、recall与预测边界数量后可知，这一差异主要来自预测边界从碎片化走向平滑时的召回变化，而不能简单解读为定位精度的好坏。
  \item \textbf{给出了跨组织的适用边界与实践建议。}在乳腺癌6名患者中，空间邻域特征相对非空间基线在每名患者中都提高了ARI，但半数患者的绝对表现很低；与BayesSpace的比较随指标和聚类数规则而变。基于这些结果，本研究对邻域层数选择、聚类数报告方式和边界评价方式提出了可操作的建议。
\end{enumerate}

\subsection{本专题的组织}
第\ref{sec:cc-background-p1}节介绍空间域识别方法、CellCharter及其已有基准、聚类数选择和空间聚类评价的相关研究。第\ref{sec:cc-methods-p1}节说明数据、评价设计与分析方法。第\ref{sec:cc-results-p1}节按上述四个问题呈现结果，并讨论其解释、实践意义与局限。第\ref{sec:cc-conclusion-p1}节总结全文。实验记录编号、数据与代码可用性说明见附录。
